\section*{Research Questions (Will be deleted in final draft)}

See my research question notes on GitHub: \url{https://github.com/Lewis-Adkins/SWLG/blob/main/SEPTransformerZoo/Research\%20Questions.md}.

\textbf{RQ1}: How does the forecasting performance of deep-learning models compare with statistical
and other machine-learning approaches in predicting high-energy solar proton flux?
This replicates Torres et al.\ and compares new transformer models against their RNN and NN.
Inputs are the same: a 2-hour window of SEP observations used to predict $t+6$ and $t+12$.
The goal is to establish a baseline of operating conditions and determine whether these models are comparable
to Torres et al.\ and to other physics-based and statistical models, such as linear regression and Posner.
The number of parameters for each model was recorded so that all models were of comparable size;
these are relatively small models.

\textbf{RQ2}: What is the relationship between forecasting performance and the scale of machine-learning models?
Having obtained an initial reading of how transformers and other models predict protons, we are interested in
adding complexity to our models by increasing the embedding vector length,
\texttt{n\_layers}, \texttt{n\_encoding\_layers}, \texttt{dim\_vals}, \texttt{batch\_sizes}, etc.
We then compare these new high-energy proton predictions to those from RQ1.
Are improvements or degradations apparent? Do all models improve at the same rate, and across all metrics?
If skill plateaus, where does that happen and why? Does the added model complexity justify the improvement?
We need to define ``acceptable improvement.''

\textbf{RQ3}: To what extent does the duration of the historical SEP observational window influence model forecasting performance?
After choosing new scales for each model at the point where skill score plateaus in RQ2, we look for further improvements.
So far, models have been trained on only two hours of data. We increase the time window to give the models more context
on the behavior of SEP electrons and protons before making a prediction.
Should we still evaluate only $t+6$ and $t+12$? How much additional training time is required?
Is the change in skill score worth the added context?

\textbf{RQ4}: At what lead time do the models' high-energy proton predictions cross a degradation threshold?
In the ML community this is referred to as ``length generalization.'' Given a fixed 2-hour SEP window,
how far into the future can we predict high-energy protons with acceptable accuracy?
The key question is how to define the ``threshold.'' There are two approaches: the first is statistical,
using residual analysis; the second, more relevant for CY27, incorporates satellite operator response time.
Where does each model cross the threshold?
